\section{A proved frozen-link reduction and its physical matching limits}
\label{nr4-section}

This section proves a nonrelativistic spectral comparison for arbitrary
\emph{fixed} links, identifies the kinetic and magnetic coefficients of
its target, and states explicit scale requirements for a proposed atomic
application. It also constructs a bounded magnetic clover modification
which preserves the determinant argument under an explicit hypothesis.
The frozen-link theorem is not a reduction of the quantized-link
Hamiltonian in a nuclear Gauss sector. In particular, its reference
vacuum uses exactly the same frozen links as its charged state.

\subsection{The exact static spectral theorem}
\label{nr4-static}

Use the lattice operators of \eqref{lat-hopping}--\eqref{lat-wilson-mass},
with $m>0$. At a fixed configuration $U$, put
\[
 R=M(U)-mI\ge0,\qquad S=\sum_{j=1}^3\sigma_jv_j(U),\qquad
 K(U)=R+\frac{S^2}{2m}.
\]
These act on the spatial two-spinor space of dimension $2N$.
In the standard Dirac representation,
\begin{equation}
 h_W(U)=\begin{pmatrix}m+R&S\\S&-m-R\end{pmatrix}.
 \label{nr4-block}
\end{equation}
No commutation of $R$ and $S$ is assumed. Write $w=\|R\|$ and
$s=\|S\|$; the uniform bounds $w\le6r/a$, $s\le3/a$ will suffice.

\begin{theorem}[Frozen-link eigenvalues, with a low-energy refinement]
\label{nr4-static-theorem}
The positive eigenvalues of \eqref{nr4-block} are $m+e_j$, in increasing
order with multiplicity, $j=1,\ldots,2N$, and its negative eigenvalues
are their opposites. If $k_j$ are the ordered eigenvalues of $K(U)$, then
\begin{equation}
 0\le e_j\le k_j,\qquad
 0\le k_j-e_j\le
 \min\left\{\frac{s^2w}{2m^2}+\frac{s^4}{8m^3},\,
                   \frac{k_j(w+k_j)}{2m}\right\}.
 \label{nr4-static-error}
\end{equation}
The first bound is uniform over the whole band and satisfies
\begin{equation}
 \delta:=\frac{s^2w}{2m^2}+\frac{s^4}{8m^3}
 \le\frac{27r}{m^2a^3}+\frac{81}{8m^3a^4}.
 \label{nr4-delta}
\end{equation}
The second bound controls the relative error on a specified low-energy
part of the spectrum. Both are independent of the number of sites.
\end{theorem}
\begin{proof}
For $-m<\lambda<m$, the lower diagonal block of $h_W-\lambda$ is
strictly negative and invertible. Its Schur complement is
\[
 m+R-\lambda+S(m+R+\lambda)^{-1}S>0.
\]
Thus there is no spectrum in $(-m,m)$. The matrix $\tau_y$ on the two
Dirac blocks anticommutes with \eqref{nr4-block}; hence the spectrum
has $2N$ opposite pairs. Eliminating the lower spinor at
$\lambda=m+e$, $e\ge0$, gives the exact equation
\begin{equation}
 A(e)u=eu,\qquad A(e)=R+S(2m+R+e)^{-1}S.
 \label{nr4-schur}
\end{equation}
Each ordered eigenvalue $a_j(e)$ of $A(e)$ is continuous and
nonincreasing. Therefore $a_j(e)-e$ is strictly decreasing, is
nonnegative at zero, and is negative for large $e$. Its unique zero
is $e_j$; this identification, including multiplicity, also follows
from Schur-complement inertia. Since $0\le A(e)\le K$, the roots
satisfy $0\le e_j\le k_j\le w+s^2/(2m)=:Q$.

For $0\le e\le Q$, functional calculus gives
\[
 0\le K-A(e)
 =S\left[(2m)^{-1}-(2m+R+e)^{-1}\right]S
 \le\frac{s^2(w+Q)}{4m^2}I=\delta I.
\]
Eigenvalue monotonicity proves the first error estimate. For the
refinement, the simpler lower bound on the resolvent yields
\[
 A(e)\ge R+\frac{S^2}{2m+w+e}
 \ge\frac{2m}{2m+w+e}K.
\]
At the $j$th root this implies
$e_j\ge2mk_j/(2m+w+e_j)$, and hence
\[
 k_j-e_j\le\frac{e_j(w+e_j)}{2m}
             \le\frac{k_j(w+k_j)}{2m}.
\]
No simultaneous diagonalization has been used.
\end{proof}

\begin{corollary}[Exact vacuum subtraction at fixed links]
\label{nr4-static-charge}
Let $E_q^{\rm stat}(U)$ be the bottom of $d\Gamma(h_W(U))$ in the
sector $Q=q$, and let $E_0^{\rm stat}(U)$ be its filled-negative-band
energy. For $|q|\le2N$,
\begin{equation}
 E_q^{\rm stat}(U)-E_0^{\rm stat}(U)
 =|q|m+\sum_{j=1}^{|q|}e_j.
 \label{nr4-static-occupation}
\end{equation}
In particular,
\begin{equation}
 0\le |q|m+\sum_{j=1}^{|q|}k_j
        -[E_q^{\rm stat}(U)-E_0^{\rm stat}(U)]\le |q|\delta.
 \label{nr4-static-charge-error}
\end{equation}
If $k_{|q|}\le E_*$, the upper bound improves to
\begin{equation}
 \frac{w+E_*}{2m}\sum_{j=1}^{|q|}k_j.
 \label{nr4-low-window}
\end{equation}
For $q=0$ both sides of the comparison vanish. A nontrivial neutral
particle--hole excitation in this quadratic static problem costs at
least $2m$.
\end{corollary}
\begin{proof}
Filling the negative levels minimizes the unconstrained quadratic
Hamiltonian. Adding $q>0$ fills its lowest $q$ positive levels;
removing $|q|$ particles creates the $|q|$ cheapest holes. Spectral
pairing gives \eqref{nr4-static-occupation}; adding further pairs only
increases the energy. Sum \eqref{nr4-static-error}. A neutral excitation
contains at least one positive particle and one negative-band hole,
each costing at least $m$.
\end{proof}

The exact lattice Pauli term retained by this theorem is
\begin{equation}
 K(U)=R+\frac1{2m}\left\{\sum_jv_j^2+
 i\sum_{j<k}\epsilon_{jk\ell}\sigma_\ell[v_j,v_k]\right\}.
 \label{nr4-pauli}
\end{equation}
This follows by multiplying Pauli matrices. The entire expression
inside braces is $S^2\ge0$, even for rough links. The leading Wilson
Laplacian $R$ cannot be omitted.

\subsection{Rest, kinetic, and magnetic matching}
\label{nr4-matching}

At $U=1$ the exact dispersion is
\begin{equation}
 E(p)=\sqrt{[m+R(p)]^2+|v(p)|^2},\quad
 R(p)=\frac r a\sum_j(1-\cos ap_j),\quad
 v_j(p)=\frac{\sin ap_j}{a}.
 \label{nr4-dispersion}
\end{equation}
At fixed $m,a,r$, its Taylor coefficient at $p=0$ gives
\begin{equation}
 E(p)=m+\frac{|p|^2}{2m_{\rm kin}}+O(|p|^4),\qquad
 \frac1{m_{\rm kin}}=\frac1m+ra,\qquad
 m_{\rm kin}=\frac{m}{1+rma}.
 \label{nr4-kinetic-mass}
\end{equation}
These are exact coefficients of this continuous-time Hamiltonian.
They are not masses imported from a spacetime-discretized lattice action.
The distinction between rest, kinetic, and magnetic masses, and the
need to match their coefficients separately, is treated in
\cite[Eq.~(9.1) and Sec.~9]{EKM}. The Hamiltonian discretization itself
is consistent with \cite[Eq.~(3)]{Zache}.

An exact free remainder, useful without a small-$pa$ assumption, is
\begin{equation}
 0\le R(p)+\frac{|v(p)|^2}{2m}-[E(p)-m]
 \le\frac{|v(p)|^2R(p)}{2m^2}+
       \frac{|v(p)|^4}{8m^3}.
 \label{nr4-free-remainder}
\end{equation}
Indeed, with $M_p=m+R(p)$,
\[
 M_p+\frac{|v|^2}{2M_p}-E
 =\frac{|v|^4}{2M_p(E+M_p)^2},
\]
and subtracting the two inverse-mass coefficients proves the bound.

We next specify the sense in which the magnetic coefficient is matched.
Take a smooth periodic classical vector potential $A$ and use the electron
convention
\[
 U_{xj}=\exp\left(-ig\int_x^{x+ae_j}A\cdot d\ell\right),\qquad
 \Pi_j=-i\partial_j+gA_j,\quad B=\nabla\times A.
\]
The same gauge-covariant difference operators can be defined on smooth
continuum functions before sampling them at lattice sites:
$T_j=e^{ia\Pi_j}$, $v_j=\sin(a\Pi_j)/a$ and
$R=(r/a)\sum_j[1-\cos(a\Pi_j)]$. These identities are just parallel
transport along a coordinate segment. Spectral calculus proves
\begin{align}
 \|(v_j-\Pi_j)f\|&\le\frac{a^2}{6}\|\Pi_j^3f\|,\nonumber\\
 \left\|\left(R-\frac{ra}{2}\sum_j\Pi_j^2\right)f\right\|
 &\le\frac{ra^3}{24}\sum_j\|\Pi_j^4f\|.
 \label{nr4-consistency}
\end{align}
These estimates are continuum $L^2$ consistency estimates on the
displayed domains, not a claim that point sampling is an isometry or
that spectra of the continuum and finite lattice have already been
compared.

For completeness, the error after squaring the spin derivative is also
controlled on smooth periodic functions. For $0<a\le a_0$, define the finite
quantity
\[
 C_{A,f,a_0}=\frac16\sum_{j,k}
 \left\{\sup_{|t|\le a_0}\|\Pi_j^3e^{it\Pi_k}\Pi_kf\|
       +\sup_{|t|\le a_0}\|\Pi_j e^{it\Pi_k}\Pi_k^3f\|\right\}.
\]
Then
\begin{equation}
 \|(S_a^2-(\sigma\cdot\Pi)^2)f\|\le a^2C_{A,f,a_0},
 \qquad S_a=\sum_j\sigma_j\frac{\sin(a\Pi_j)}a.
 \label{nr4-squared-consistency}
\end{equation}
To prove it, expand each difference as
$(v_j-\Pi_j)v_k+\Pi_j(v_k-\Pi_k)$, apply
\eqref{nr4-consistency}, and use
$v_kf=(2a)^{-1}\int_{-a}^ae^{it\Pi_k}\Pi_kf\,dt$.
The second summand is bounded by applying the integral Taylor remainder
for $v_k-\Pi_k$ before $\Pi_j$. In detail, the average of the linear
term in $e^{it\Pi_k}$ vanishes, and its quadratic integral remainder
has total absolute weight $a^2/6$. The supremum in the second term
of $C_{A,f,a_0}$ is needed because $\Pi_j$ and $\Pi_k$ need not commute.

Since $[\Pi_j,\Pi_k]=-ig\epsilon_{jk\ell}B_\ell$,
$(\sigma\cdot\Pi)^2=\Pi^2+g\sigma\cdot B$. Consequently the leading
smooth-field consistency operator is
\begin{equation}
 K_{\rm cont}^{\rm tree}
 =\frac{\Pi^2}{2m_{\rm kin}}+\frac{g}{2m}\sigma\cdot B.
 \label{nr4-tree-pauli}
\end{equation}
The norm difference on $f$ is bounded by the second line of
\eqref{nr4-consistency} plus $a^2C_{A,f,a_0}/(2m)$.
Defining the tree-level spin factor through the coefficient
$g\,g_{\rm eff}\sigma\cdot B/(4m_{\rm kin})$ gives
\begin{equation}
 g_{\rm eff}=\frac{2m_{\rm kin}}m=\frac2{1+rma}.
 \label{nr4-g-factor}
\end{equation}
This is coefficient matching for smooth weak fields and low momenta.
It is not a calculation of the interacting electron's measured magnetic
moment, which also requires radiative and parameter matching.

\subsection{An explicit hierarchy with a scaled Wilson parameter}
\label{nr4-scaled}

Set $x=ma$ and choose
\begin{equation}
 r=\frac c x,\qquad c>0.
 \label{nr4-scaled-r}
\end{equation}
Then $m_{\rm kin}=m/(1+c)$ and $g_{\rm eff}=2/(1+c)$.
The fractional spin-factor error from $2$ is $c/(1+c)$, and its
absolute error is $2c/(1+c)$. At a free Brillouin-zone corner with
$k$ components equal to $\pi/a$, the energy is exactly $m+2kr/a$.
The first such excess threshold is therefore
\begin{equation}
 \Delta_d=\frac{2r}a=\frac{2cm}{x^2}.
 \label{nr4-doubler-scale}
\end{equation}
Corners occur exactly for even periodic lattices; these formulas also
describe the large-volume band structure. They identify free thresholds,
not a proof that rough gauge fields or nuclear potentials cannot mix
the corresponding bands.
A conservative free spectral exclusion is immediate: if any
$|ap_j|\ge\pi/2$ in the first Brillouin zone, then $R(p)\ge r/a$ and
$E(p)-m\ge r/a=\Delta_d/2$. Thus a free energy cutoff strictly below
$\Delta_d/2$ excludes the outer half of every momentum coordinate.
Calling every corner a local minimum would require a further parameter
restriction; the stated corner energies and this exclusion do not.

For a stated hydrogenic benchmark of nuclear charge $Z$, define
\begin{equation}
 p_Z=Z\alpha m_{\rm kin},\qquad
 \mathcal R_Z=\frac{(Z\alpha)^2m_{\rm kin}}2.
 \label{nr4-hydrogenic-scales}
\end{equation}
These are the momentum and binding-energy scales of the continuum
Coulomb problem with kinetic mass $m_{\rm kin}$. Their use here is a
parameter test, not an estimate on an interacting lattice ground state.

\begin{proposition}[Quantified compatibility of the benchmark scales]
\label{nr4-hierarchy}
Fix $H>1$, momentum-resolution tolerance $0<\varepsilon_a<1$, and
energy-separation tolerance $0<\varepsilon_d<1$. The three requirements
\[
 x\ge H,\qquad ap_Z\le\varepsilon_a,\qquad
 \mathcal R_Z/\Delta_d\le\varepsilon_d
\]
hold precisely when
\begin{equation}
 H\le x\le\min\left\{
 \frac{\varepsilon_a(1+c)}{Z\alpha},\,
 \frac{2\sqrt{\varepsilon_dc(1+c)}}{Z\alpha}\right\}.
 \label{nr4-window}
\end{equation}
If the frozen-link target levels additionally satisfy $k_j\le\mathcal R_Z$,
their relative spectral error from Theorem~\ref{nr4-static-theorem} is
bounded by
\begin{equation}
 \frac{k_j-e_j}{k_j}\le\frac{3c}{x^2}
                      +\frac{(Z\alpha)^2}{4(1+c)}
 \quad (k_j>0).
 \label{nr4-scaled-relative-error}
\end{equation}
The error vanishes exactly when $k_j=0$.
\end{proposition}
\begin{proof}
Direct substitution gives $ap_Z=Z\alpha x/(1+c)$ and
$\mathcal R_Z/\Delta_d=(Z\alpha)^2x^2/[4c(1+c)]$, proving the
equivalence. Also $w/m\le6c/x^2$; insert this and
$\mathcal R_Z/m=(Z\alpha)^2/[2(1+c)]$ in
\eqref{nr4-static-error}.
\end{proof}

At $c=0.1$ and $\alpha=1/137$, the spin factor is $20/11$ and its
fractional error is $1/11$, about $9.09\%$. For $Z=1$, $x=60$ gives
$ap_Z\simeq0.398$ and $\mathcal R_Z/\Delta_d\simeq0.436$.
Thus ``about sixty'' is at best an order-of-magnitude crossover;
it does not express two strong scale separations. With the explicit
choices $\varepsilon_a=\varepsilon_d=0.1$, the upper limit in
\eqref{nr4-window} is $15.07$ for $Z=1$.

For the unscreened $Z=26$ benchmark relevant to inner iron electrons,
even replacing both tolerances by $1$ gives an upper limit about $3.50$.
With the $10\%$ choices it becomes about $0.580$, incompatible even
with $x>1$. Valence scales with a smaller effective charge do not
resolve this inner-shell requirement for a theorem uniform up to
$Z=26$. This calculation does not prove that an improved regulator
cannot describe iron. It shows that the proposed simultaneous
heavy-mass, small-$c$, well-resolved single-species hierarchy cannot
be justified for that benchmark merely by citing the hydrogen value.

\subsection{A magnetic clover term: its actual symmetry and a sufficient bound}
\label{nr4-clover}

In the formulas for magnetic field normalization, choose
$g=+\sqrt{g^2}>0$.
For each cyclic triple $(j,k,\ell)$, let the four plaquettes in the
$jk$ plane adjacent to $x$ all have orientation $(j,k)$. Define a real
Abelian magnetic clover multiplier by
\begin{equation}
 B_\ell^{\rm cl}(x;U)
 =-\frac1{4ga^2}\sum_{p\in\mathcal C_{jk}(x)}\operatorname{Im}U_p.
 \label{nr4-clover-definition}
\end{equation}
For smooth links with the convention above this is $B_\ell(x)+O(a^2)$;
oppositely placed plaquettes cancel first derivative errors. It is
gauge invariant for $U(1)$ and is a bounded continuous function on
the full compact link space. In particular,
\begin{equation}
 |B_\ell^{\rm cl}(x;U)|\le\frac1{ga^2},\qquad
 \|\Sigma\cdot B^{\rm cl}(U)\|\le\frac{\sqrt3}{ga^2},
 \label{nr4-clover-norm}
\end{equation}
where $\Sigma_\ell=-i\alpha_j\alpha_k$ for cyclic $(j,k,\ell)$.
The four-leaf field-strength construction is standard; see
\cite[Eq.~(2.12)]{EKM}. Our real-field normalization and continuous-time
magnetic Hamiltonian are specified independently by
\eqref{nr4-clover-definition}.

Consider the explicitly different one-particle Hamiltonian
\begin{equation}
 h_\lambda(U)=h_W(U)+\lambda\beta\Sigma\cdot B^{\rm cl}(U),
 \qquad \lambda\in\mathbb R.
 \label{nr4-clover-h}
\end{equation}
The often relevant symmetry test can be made without any convention
about the name of a chirality matrix. Our $C_5=\alpha_1\alpha_2\alpha_3\beta$
anticommutes with each $\alpha_j$ and with $\beta$, and therefore
\begin{equation}
 [C_5,\Sigma_\ell]=0,\qquad
 \{C_5,\beta\Sigma_\ell\}=0,\qquad
 \tfrac12\{\beta,h_\lambda\}
 =M(U)+\lambda\Sigma\cdot B^{\rm cl}(U).
 \label{nr4-clover-algebra}
\end{equation}
Thus the clover term \emph{preserves} the $C_5$ anticommutation.
It changes the positive operator needed in the $\beta$-form argument;
its contribution is not positive by itself.

\begin{proposition}[A compact-link clover variant with vacuum stability]
\label{nr4-clover-positive}
Replace $h_W$ by $h_\lambda$ in the full compact-link Hamiltonian and
retain the same electric energy, magnetic plaquette energy, and local
Gauss operators. If
\begin{equation}
 m_{\rm cl}:=m-\frac{\sqrt3|\lambda|}{ga^2}\ge0,
 \label{nr4-clover-mass}
\end{equation}
then the physical-sector comparison remains valid, with
\begin{equation}
 E_\nu^{(\lambda)}\ge E_0^{(\lambda)}
                         +m_{\rm cl}\left|\sum_x\nu_x\right|.
 \label{nr4-clover-stability}
\end{equation}
The same prescribed nuclear self-energy may be subtracted; the
linear nuclear bound then follows with $m_{\rm cl}$ in its rest-cost
term. No restriction on individual link configurations is imposed.
\end{proposition}
\begin{proof}
The new term is bounded, Hermitian, gauge covariant, and traceless in
the full one-particle space. It preserves anticommutation with $C_5$.
By \eqref{nr4-clover-norm},
$\{\beta,h_\lambda(U)\}/2\ge m_{\rm cl}I$ uniformly in $U$.

The determinant and annulus proofs need only these algebraic
properties: $\{C_5,h\}=0$, $\{\beta,h\}/2\ge m_{\rm cl}I$,
and scalar kicks commuting with $C_5$ and $\beta$.
Indeed $U^*C_5U=C_5$ still proves centered-trace reality, while
\[
 \frac d{dt}\langle f,\beta f\rangle
 =-\langle f,\{\beta,h_\lambda\}f\rangle
 \le-2m_{\rm cl}\|f\|^2
\]
proves nonvanishing and the quantitative annulus when $m_{\rm cl}>0$.
For positivity of the real nonzero trace, deform $h_\lambda(t)$ linearly
to $m_{\rm cl}\beta$ and turn off the scalar kicks. This preserves
both anticommutation and the displayed lower bound. The terminal
trace is positive. Thus the bounded matrix Feynman--Kac and full local
Gauss projection proof, including its global-charge Fourier estimate,
applies unchanged. A bounded mass perturbation followed by a limit
covers $m_{\rm cl}=0$.
\end{proof}

In smooth weak fields the new Pauli coefficient is
$g/(2m)+\lambda$. This coefficient can also be characterized by a
finite-matrix small-parameter statement: if the upper mass correction
is $\varepsilon^2D+O(\varepsilon^4)$ and the off-diagonal block is
$\varepsilon S_1+O(\varepsilon^3)$ in operator norm, with fixed bounded
Hermitian $D,S_1$, the positive energies are
\[
 m+\varepsilon^2\lambda_j\left(D+\frac{S_1^2}{2m}\right)
       +O(\varepsilon^4).
\]
Indeed Weyl's inequality first puts these energies within $O(\varepsilon)$
of $m$. The lower-block Schur complement then gives $e=O(\varepsilon^2)$,
and its inverse differs from $(2m)^{-1}I$ by $O(\varepsilon^2)$.
Substitution and eigenvalue monotonicity prove the displayed expansion.
This argument allows $D$ to be indefinite; it applies to the coefficients
where covariant momentum has order $\varepsilon$ and magnetic field
has order $\varepsilon^2$. It identifies coefficients and does not assert
that an arbitrary quantum link state obeys this scaling.
Consequently choosing
\begin{equation}
 \lambda_{\rm tree}=\frac{gra}{2}
 \label{nr4-clover-tree}
\end{equation}
makes the tree-level spin factor equal to $2$ with the kinetic mass
in \eqref{nr4-kinetic-mass}. The sufficient positivity requirement is
then simply
\begin{equation}
 m\ge\frac{\sqrt3r}{2a}.
 \label{nr4-clover-tree-positive}
\end{equation}
For $r=c/x$ it becomes $x^2\ge\sqrt3c/2$, and
$m_{\rm cl}=m[1-\sqrt3c/(2x^2)]$.
Thus a bounded magnetic clover term does not automatically destroy the
vacuum theorem in the heavy regime. This conclusion uses the compact
field bound and the specified coefficient; it is not an assertion
that an arbitrary continuum Pauli term is harmless. The modification
does not by itself match electric, Darwin, spin--orbit, or radiative
terms, or change the free doubler thresholds, since it vanishes at
$U=1$. A term involving electric rotor operators requires a separate
domain and positivity analysis. Distinguishing magnetic and electric
improvement is also necessary in the transfer-matrix treatment of
\cite[Sec.~5.1]{EKM}.

\subsection{What the static theorem does not transfer to quantized links}
\label{nr4-remaining}

For fixed $U$, the sea energy in the exact comparison is
\[
 E_0^{\rm stat}(U)=-\sum_{j=1}^{2N}[m+e_j(U)].
\]
When $U$ is quantized, this is a multiplication operator depending on
the links. The interacting source state and vacuum need not have the
same expectation of it. It cannot be discarded as a common scalar.
Moreover, even on a chart where a smooth band-diagonalizing unitary
$\mathcal U(U)$ is available, transformation of the electric derivative
produces
\[
 \mathcal U E_\ell\mathcal U^*
 =E_\ell+i(\partial_{\varphi_\ell}\mathcal U)\mathcal U^*.
\]
The induced connection and its off-band parts are additional terms in
the electric energy. A comparison of frozen eigenvalues does not bound
them. This elementary identity already explains why substituting
$R+S^2/(2m)$ into a quantized Pauli--Fierz-type Hamiltonian is not a
proof of that reduction.

The frozen-link charge-extensive error is therefore established only
for \eqref{nr4-static-occupation}, with a common $U$ and no scalar nuclear
potential or Gauss constraint. A full low-energy theorem still needs
a specified retained subspace, a controlled elimination of its
complement with the electric dynamics included, the induced gauge and
vacuum-polarization terms, and an error relative to the actual source-free
vacuum. Neither the existing charge gap nor the static $2m$ pair cost
supplies these missing estimates for the quantized theory.
